\documentclass[border=4pt]{standalone}
% Shared CircuiTikZ preamble for the Lorenz analog circuit diagrams.
\usepackage{tikz}
\usepackage{amsmath}
\usepackage{circuitikz}
\ctikzset{
  amplifiers/fill=white,
  resistors/scale=0.9,
  capacitors/scale=0.9,
  bipoles/length=0.9cm,
}


\begin{document}
\begin{circuitikz}[american]

  % AD633 as a labeled rectangle with named pin coordinates.
  \def\chw{2.6}   % chip width
  \def\chh{3.2}   % chip height
  \coordinate (chipNW) at (0,\chh);
  \coordinate (chipSE) at (\chw,0);
  \draw[thick] (0,0) rectangle (\chw,\chh);
  \node at ($(chipNW)!0.5!(chipSE)$) [font=\bfseries\small] {AD633};
  \node at ($(0,0)!0.5!(\chw,0)+(0,-0.9)$) [font=\scriptsize, align=center]
      {$W = \dfrac{(X_1-X_2)(Y_1-Y_2)}{10\,\mathrm{V}} + B$};

  % Input pin coordinates (left edge).
  \coordinate (X1) at (0, \chh-0.4);
  \coordinate (X2) at (0, \chh-1.0);
  \coordinate (Y1) at (0, \chh-1.8);
  \coordinate (Y2) at (0, \chh-2.4);
  \coordinate (Bin) at (0, \chh-3.0);

  % Pin labels just inside the chip.
  \node at ($(X1)+(0.25,0)$)  [right, font=\scriptsize] {$X_1$};
  \node at ($(X2)+(0.25,0)$)  [right, font=\scriptsize] {$X_2$};
  \node at ($(Y1)+(0.25,0)$)  [right, font=\scriptsize] {$Y_1$};
  \node at ($(Y2)+(0.25,0)$)  [right, font=\scriptsize] {$Y_2$};
  \node at ($(Bin)+(0.25,0)$) [right, font=\scriptsize] {$B$};

  % Output pin (right edge).
  \coordinate (W) at (\chw, \chh-1.5);
  \node at ($(W)+(-0.25,0)$) [left, font=\scriptsize] {$W$};

  % Wiring to produce W = -XY/10V (the Y stage supplies -Y):
  %   X1 = X, X2 = ground   (so X1-X2 = X)
  %   Y1 = -Y, Y2 = ground  (so Y1-Y2 = -Y)
  %   B (summing input) grounded, so B = 0
  %   --> W = X * (-Y) / 10V + 0 = -XY/10V
  \draw (X1) -- ++(-1.6,0) node[circ]{} node[left] {$X$};
  \draw (X2) -- ++(-0.8,0) node[ground, rotate=-90, yshift=-0.05cm]{};
  \draw (Y1) -- ++(-1.6,0) node[circ]{} node[left] {$-Y$};
  \draw (Y2) -- ++(-0.8,0) node[ground, rotate=-90, yshift=-0.05cm]{};
  \draw (Bin) -- ++(-0.8,0) node[ground, rotate=-90, yshift=-0.05cm]{};

  % Output wire.
  \draw (W) -- ++(1.6,0) node[circ]{} node[right] {$-\dfrac{XY}{10\,\mathrm{V}}$};

\end{circuitikz}
\end{document}
